\chapter{Mucus in Urine}
\section*{What Is This Test?} Mucus in urine is assessed during urinalysis microscopy to determine whether threads or strands are present and clinically meaningful.\section*{Alternative Names} Urine sediment mucus; urinary mucus threads.\section*{Principle} Microscopy identifies mucus and other sediment elements; small amounts can be normal because the urinary tract produces mucus.\section*{Why This Test Is Ordered} It may be evaluated with urinary pain, discharge, infection symptoms, stones, or an abnormal urinalysis.\section*{Primary vs Secondary Test} Urinalysis is primary; culture, STI testing, or imaging is selected based on symptoms.\section*{How This Test Is Performed} A clean-catch urine sample is analyzed by dipstick and microscopy.\section*{What Preparation Is Needed} Follow clean-catch instructions and report menstruation or genital discharge that could contaminate the sample.\section*{Understanding Results} Small amounts are often nonspecific; abundant mucus with white cells, bacteria, blood, or symptoms may warrant evaluation.\section*{Follow-up Tests} Urine culture, STI testing, repeat urinalysis, or imaging may follow.\section*{References} \begin{enumerate}\item MedlinePlus. Mucus in Urine.\item American Urological Association. Urinalysis resources.\end{enumerate}\clearpage\section*{Understanding Results and Follow-up}
The laboratory report should identify the specimen, method, units, reference interval or decision limit, and whether the result is positive, negative, abnormal, or inconclusive. Reference intervals vary with age, sex, pregnancy, specimen type, assay, and laboratory; a result outside a range is not automatically a diagnosis, and a result within a range does not always exclude disease. Discuss unexpected results with the ordering clinician, who can decide whether repeat or confirmatory testing, treatment, or referral is appropriate. Do not change medicines or delay urgent care based on an isolated result.
\section*{Regulatory Status and Authorization Date}This chapter describes a test category, not a specific commercial product. FDA, CDSCO, EU IVDR, and MHRA approval or registration dates therefore cannot be assigned without the assay manufacturer, product name, instrument, and intended use. For laboratory-developed tests, an individual agency approval date may not exist. See the Preface for the interpretation of regulatory status.
\clearpage
